\section{Core Algorithms}
\label{sec:algorithms}

\subsection{StringToDAG (S2D) --- The Decoder}

The S2D decoder executes an IsalSR instruction string token-by-token to
build a \texttt{LabeledDAG}. It tokenizes the input using a greedy scanner
(consuming two characters for $V/v$ compound tokens) and dispatches each
token to update the shared state: (DAG, CDLL, primary pointer, secondary pointer).

\begin{algorithm}[h]
\caption{StringToDAG (S2D)}
\begin{algorithmic}[1]
\Require String $w$, number of variables $m$
\State Initialize DAG with $m$ VAR nodes; CDLL $\leftarrow [x_1, \ldots, x_m]$
\State $\text{pri} \leftarrow \text{CDLL node for } x_1$; $\text{sec} \leftarrow \text{pri}$
\For{each token $t$ in $\textsc{Tokenize}(w)$}
  \If{$t \in \{N, P, n, p\}$} move corresponding pointer
  \ElsIf{$t = V\ell$} create node with label $\ell$, edge pri$\to$new, insert after pri
  \ElsIf{$t = v\ell$} same from secondary pointer
  \ElsIf{$t = C$} add edge pri$\to$sec if acyclic
  \ElsIf{$t = c$} add edge sec$\to$pri if acyclic
  \EndIf
\EndFor
\State \Return DAG
\end{algorithmic}
\end{algorithm}

\subsection{DAGToString (D2S) --- The Greedy Encoder}

The D2S encoder converts a \texttt{LabeledDAG} back into an instruction string
using a greedy algorithm adapted from IsalGraph \cite{lopezrubio2025}.
It searches for the minimum-cost pointer displacement that enables a structural
operation (V, v, C, or c) at each step.

The displacement pairs $(a, b)$ are sorted by $|a| + |b|$ (total movement cost),
ensuring the shortest possible movement instructions are emitted first.
The priority order is $V > v > C > c$: node insertion is preferred over edge-only
insertion.

Key adaptations from IsalGraph:
\begin{itemize}[nosep]
  \item Token emission: \texttt{"V" + label\_char} (two-character compound tokens).
  \item Neighbor search: \texttt{out\_neighbors()} only (directed data flow).
  \item All $m$ VAR nodes are pre-mapped at initialization.
\end{itemize}

\subsection{Round-Trip Property}

\begin{theorem}[Round-trip]
For any valid IsalSR string $w$ and number of variables $m$:
\[
  \text{S2D}(w, m) \cong \text{S2D}(\text{D2S}(\text{S2D}(w, m), x_1), m)
\]
where $\cong$ denotes labeled-DAG isomorphism.
\end{theorem}

This property is verified by 25 deterministic tests and Hypothesis property
tests across random strings.

\subsection{Canonical String --- The Core Contribution}

\begin{definition}[Canonical string]
For a labeled DAG $D$ with initial node $x_1$:
\[
  w^*_D = \text{lexmin}\left\{ w \in \arg\min_{P \in \text{Paths}(D, x_1)} |P| \right\}
\]
The shortest D2S string from $x_1$, with ties broken lexicographically.
Computed via exhaustive backtracking over all valid neighbor choices at V/v
branch points.
\end{definition}

\begin{theorem}[Complete invariant]
\label{thm:invariant}
$w^*_{D_1} = w^*_{D_2}$ if and only if $D_1 \cong D_2$ (labeled-DAG isomorphism).
\end{theorem}

\textbf{Simplification from IsalGraph}: Since input variables are distinguishable
($x_i$ maps to $x_i$ under any isomorphism), the canonical string is computed
from $x_1$ only---no iteration over starting nodes. This reduces the canonical
computation cost by a factor of $O(n)$.

\subsection{Pruned Canonical with 6-Component Structural Tuples}

At each V/v branch point, only candidates sharing the maximum 6-component
structural tuple are explored:
\[
  \tau(v) = (|\text{in}_{N_1}(v)|, |\text{out}_{N_1}(v)|,
             |\text{in}_{N_2}(v)|, |\text{out}_{N_2}(v)|,
             |\text{in}_{N_3}(v)|, |\text{out}_{N_3}(v)|)
\]
where $\text{in}_{N_k}(v)$ and $\text{out}_{N_k}(v)$ count nodes at distance $k$
following incoming and outgoing edges, respectively. This is computed via
truncated BFS in both directions.

The pruned variant preserves the complete invariant property while dramatically
reducing the branching factor (empirically 10--100$\times$ faster).

\subsection{DAG Distance Metric}

The Levenshtein distance between canonical strings defines a metric on
labeled DAGs:
\[
  d(D_1, D_2) = \text{Levenshtein}(w^*_{D_1}, w^*_{D_2})
\]
This satisfies symmetry, triangle inequality, and identity of indiscernibles.
